\scanpage{107}
% archive foliation: 106
\textcircled{2} \uline{Cas \deleted{\unsure{non}} lisse non quasi-projectif.}

Si $X$ est séparé, utilisant une compactification de Nagata, et résolution des singularités \inserted{comme dessus}, on a les \unsure{ms} résultats que ci-dessus, en notant que la résolution des singularités permet de déduire les conj. de Weil pour des var. propres et lisses du cas projectif et lisse.

Si $X$ non séparé, $H^i_!(X)$ n'est plus défini. Par contre $H^i(X)$ est défini. On a, avec un \unsure{recouvrt} affine
\[
H^{\unsure{*}}(X)\Longleftarrow E^{pq}_2=H^p(\sigma\longleftarrow H^q(U_\sigma))
\]
Ici $E^{pq}_2$ est effectif de poids $\unsure{\ge|\not\ge}q$, mais \deleted{\unsure{malgré}} \deleted{\unsure{faible degré}} on peut en conclure que l'effectivité de $H^i(X)$, et le fait que ses poids sont compris entre $0$ et $2n$ si $i\ge n$, et entre $0$ et $2i$ si $i\le n$; en tous cas, ces poids sont $\le2i$, \deleted{\unsure{enfin le ni}} Enfin, on trouve que le niveau est $\le i$ \deleted{\unsure{en tous cas}, \unsure{et $\le2n-i$} (il y a de \unsure{singuliers} pour $i>n$)}.

\margin{\fbox{\parbox{.26\linewidth}{\uline{Question} le niveau ne serait-il $\le2n-i$ ? \deleted{\unsure{pour $i\ge n$}}}}}
\margin{$n=\dim X$}

\textcircled{3} \uline{Cas non lisse} : $H^i_!(X)$, ($X$ séparé). Je dis qu'il est encore vrai qu'il est \uline{effectif} et de poids $\le i$, de niveau $\operatorname{Inf}(i,2n-i)$.
Soit $U$ un ouvert lisse dense, $Y$ son complémentaire.
Donc
\diagram{\begin{tikzcd}
{} \arrow[r] & H^i_!(U)\deleted{\unsure{f}} \arrow[r] & H^i_!(X) \arrow[r] & H^i_!(Y) \arrow[r] & {}
\end{tikzcd}}
et $H^i_!(U)$ \deleted{\unsure{est}} est effectif de poids $\le i$ \inserted{de niveau $\le i$}, et $H^i_!(Y)$ aussi par hypothèse de récurrence sur la dimension, donc $H^i_!(X)$ est effectif de poids $\le i$, de niveau $\le i$. De \(=\) $H^i_!(U)$ est de niveau $\le2n-i$, et $H^i_!(Y)$ de niveau $\le2(n-1)-i\le2n-i$, on \unsure{gagne}.

\textcircled{4} \uline{Cas non lisse} : $H^i(X)$ \deleted{($X$ \unsure{séparé})}. Je dis qu'il est effectif, à poids $\uuline{\le2i}$, à niveau \deleted{$\unsure{\le\operatorname{Inf}(i,2n-i)}$ et $n=\dim X$} \inserted{$\le i$ et $\le n$}.
\inserted{Par la suite spectrale de Leray, on est ramené au cas où $X$ affine.}
Récurrence sur $i$; trivial si $i\le0$. On utilise résolution $X'\to X$, et suite spectrale de descente
\[
H^{\unsure{*}}(X)\Longleftarrow E^{pq}_2=H^p(\unsure{s|\sigma}\longmapsto H^q(\unsure{(X'/X)^s})),
\]
\deleted{\unsure{en}} \deleted{\unsure{donc}} Par hyp. de récurrence, \deleted{\unsure{pour}} \inserted{si $q<i$,} $E^{pq}_2$ est effectif

\margin{poids $\le2\operatorname{Inf}(i,n)$\\niveau $\le\operatorname{Inf}(i,n)$}
